\section*{Chapter \ref{ch:b2:mantle-melting} --- Melting the Mantle}

\begin{solution}{exo:b2:mantle-melting:1}
(a) More than \qty{90}{\percent} olivine: \omterm{def:b2:mantle-melting:peridotites}{dunite}. (b) Olivine and orthopyroxene
without clinopyroxene: \omterm{def:b2:mantle-melting:peridotites}{harzburgite}. (c) More than \qty{40}{\percent} olivine with
both pyroxenes: \omterm{def:b2:mantle-melting:peridotites}{lherzolite}. (d) Less than \qty{40}{\percent} olivine: a pyroxenite
(an olivine-websterite, with both pyroxenes). (e) $57/98 = \qty{58.2}{\percent}$,
$28/98 = \qty{28.6}{\percent}$, $13/98 = \qty{13.3}{\percent}$: a \omterm{def:b2:mantle-melting:peridotites}{lherzolite}.
\end{solution}

\begin{solution}{exo:b2:mantle-melting:2}
(a) The liquidus at \qty{60}{\percent} anorthite is \qty{1390}{\celsius}, in the
anorthite field: anorthite crystallises first. (b) Diopside appears at the
eutectic, \qty{1270}{\celsius}, when the liquid has \qty{42}{\percent} anorthite.
(c) The first liquid of any mixture of the two minerals is the eutectic liquid,
\qty{42}{\percent} anorthite, at \qty{1270}{\celsius}.
\end{solution}

\begin{solution}{exo:b2:mantle-melting:3}
(a) $\num{4.0e-5}\times1623/(3300\times1000) = \qty{1.97e-8}{K/Pa} =
\qty{19.7}{K/GPa}$. (b) $\num{e9}/(3300\times9.81) = \qty{3.09e4}{m}$, or
\qty{30.9}{km} per gigapascal; $19.7/30.9 = \qty{0.64}{K/km}$, about fifty times
less than the \qty{30}{K/km} near the surface, where heat is conducted, not
carried.
\end{solution}

\begin{solution}{exo:b2:mantle-melting:4}
(a) $T_{\mathrm s}(0) = \qty{1085.7}{\celsius}$; $T_{\mathrm s}(1) = 1085.7 +
132.9 - 5.1 = \qty{1213.5}{\celsius}$; $T_{\mathrm s}(2) = 1085.7 + 265.8 - 20.4 =
\qty{1331.1}{\celsius}$; $T_{\mathrm s}(3) = 1085.7 + 398.7 - 45.9 =
\qty{1438.5}{\celsius}$. (b) $132.9 - 2\times5.1\times2 = \qty{112.5}{K/GPa}$.
(c) At \qty{3}{GPa}, \qty{1400}{\celsius} is below the solidus,
\qty{1438.5}{\celsius}; rising to \qty{2}{GPa} along the adiabat it cools by only
about \qty{20}{\celsius}, to some \qty{1380}{\celsius}, while the solidus falls to
\qty{1331.1}{\celsius}: the rock is now above its solidus and part of it melts.
\end{solution}

\begin{solution}{exo:b2:mantle-melting:5}
(a) $1350 + (1277.5 - 1350)\times(30 - 20)/(40 - 20) = \qty{1314}{\celsius}$.
(b) The composition lies on the diopside side of the eutectic: diopside.
(c) The liquid reaches the eutectic with \qty{42}{\percent} anorthite; anorthite
stays in the liquid, so the liquid fraction is $30/42 = 0.714$ and $1 - 0.714 =
\qty{29}{\percent}$ of the system has crystallised as diopside. (d) Everything
crystallises: \qty{70}{\percent} diopside and \qty{30}{\percent} anorthite, the
bulk composition.
\end{solution}

\begin{solution}{exo:b2:mantle-melting:6}
(a) Forsterite: $60.08/(2\times40.30 + 60.08) = 60.08/140.68 =
\qty{42.71}{\percent}$ \ce{SiO2}; enstatite: $60.08/100.38 = \qty{59.85}{\percent}$.
(b) $0.975\times59.85 + 2.5 = \qty{60.85}{\percent}$. (c) Enstatite lies between
forsterite and the liquid: the fraction of forsterite is $(60.85 - 59.85)/(60.85 -
42.71) = 1.00/18.14 = \qty{5.5}{\percent}$, Bowen and Andersen's value. (d) The
melt is richer in silica than the orthopyroxene it comes from, and the residue
gains olivine.
\end{solution}

\begin{solution}{exo:b2:mantle-melting:7}
(a) $5.1P^2 - 112.9P + 214.3 = 0$: $P_0 = (112.9 - \sqrt{12746.4 -
4371.7})/10.2 = (112.9 - 91.51)/10.2 = \qty{2.10}{GPa}$, at $2.097\times30.9 =
\qty{64.8}{km}$. (b) $5.1P^2 - 112.9P + 364.3 = 0$: $P_0 = (112.9 -
72.90)/10.2 = \qty{3.92}{GPa}$, at \qty{121}{km}. (c) $T_{\mathrm s}'(2.10) = 132.9
- 21.4 = \qty{111.5}{K/GPa}$, so $\Delta P_0 = 10/(111.5 - 20) = \qty{0.109}{GPa}$
per \qty{10}{\celsius}, or \qty{3.4}{km}.
\end{solution}

\begin{solution}{exo:b2:mantle-melting:8}
(a) $T_{\mathrm s} = \qty{1331.1}{\celsius}$, $T_{\mathrm l} = 1475.0 + 160.0 -
12.8 = \qty{1622.2}{\celsius}$; $T = 1331.1 + 0.10\times291.1 =
\qty{1360.2}{\celsius}$, or \qty{1633}{K}. (b) $(\partial T/\partial P)_F = 112.5
+ 0.10\times(67.2 - 112.5) = \qty{108.0}{K/GPa}$; $(\partial T/\partial F)_P =
T_{\mathrm l} - T_{\mathrm s} = \qty{291.1}{K}$. (c) $c_p/T = 0.612$; with $\alpha/\rho
= \qty{12.1}{J/(kg.K)}$ per gigapascal, the numerator is $0.612\times108.0 - 12.1
= 54.0$ and the denominator $300 + 0.612\times291.1 = 478.2$: $-\mathrm dF/\mathrm
dP = 0.113$, \qty{11.3}{\percent} per gigapascal. (d) $54.0/178.2 = 0.303$,
\qty{30}{\percent} per gigapascal: the latent heat cuts the productivity by a
factor of 2.7.
\end{solution}

\begin{solution}{exo:b2:mantle-melting:9}
(a) $43\times1 = \qty{43}{\celsius}$; $43\times5^{0.75} = 43\times3.34 =
\qty{144}{\celsius}$; $43\times13^{0.75} = 43\times6.85 = \qty{294}{\celsius}$.
(b) $1213.5 - 294 = \qty{919}{\celsius}$. (c) \qty{125}{ppm} is
\qty{0.0125}{\percent}, so $X = 0.0125/0.01 = \qty{1.25}{\percent}$ and $\Delta T =
43\times1.25^{0.75} = 43\times1.18 = \qty{51}{\celsius}$.
\end{solution}

\begin{solution}{exo:b2:mantle-melting:10}
(a) The base is $2H$ wide and the mantle crosses it at $U$: $2UH$ per unit length
of ridge and time. Each side spans the height $H$ and is crossed horizontally at
$U$: $UH$ each, $2UH$ in all. With sides at $\theta$ the base would be
$2H/\tan\theta$ wide and the inflow $2UH/\tan\theta$, equal to the outflow only if
$\tan\theta = 1$: in this model the angle is set by mass conservation. (b) With
constant productivity the \omterm{def:b2:mantle-melting:fraction}{melt fraction} at the apex is $F_{\max} = \Pi P_0 = \Pi\rho
g z_0$, and $H = z_0$: $h_c = F_{\max}z_0/2 = \Pi\rho g z_0^2/2$. (c) Each year
\qty{50}{mm} of new plate, $\qty{0.05}{m}\times\qty{1000}{m}\times\qty{7100}{m} =
\qty{3.6e5}{m^3}$ of crust per kilometre of ridge; at \qty{12}{\percent} mean \omterm{def:b2:mantle-melting:fraction}{melt
fraction} it comes from $\num{3.6e5}/0.12 = \qty{3.0e6}{m^3}$ of mantle.
\end{solution}

\begin{solution}{exo:b2:mantle-melting:11}
(a) After a mass $t$ of A has left a unit mass of liquid, the liquid holds $0.6 -
t$ of A, 0.2 of B and 0.2 of C in a mass $1 - t$: $\big((0.6 - t), 0.2,
0.2\big)/(1 - t)$, and B/C stays 1. (b) On the cotectic, $(0.40 - 0.15s, 0.60 -
0.30s, 0.45s)$; B = C gives $0.60 - 0.30s = 0.45s$, $s = 0.8$: $L_1$ has A 28, B
36, C 36. Then $0.2/(1 - t) = 0.36$, $t = 0.44$: \qty{44}{\percent} of the system has
crystallised as A. (c) C is all in the liquid: $f = 20/45 = 0.44$,
\qty{44}{\percent} still liquid. (d) A: $0.60 - 0.444\times0.25 = 0.49$; B: $0.20 -
0.444\times0.30 = 0.07$; with the liquid, $49 + 7 + 44 = \qty{100}{\percent}$.
\end{solution}

\begin{solution}{exo:b2:mantle-melting:12}
(a) $X_{\mathrm{sat}}(1) = 12 + 1 = \qty{13}{\percent}$, $\Delta T =
\qty{294}{\celsius}$, saturated solidus $1213.5 - 294 = \qty{919}{\celsius}$;
$X_{\mathrm{sat}}(2) = 12\times1.516 + 2 = \qty{20.2}{\percent}$, $\Delta T =
43\times20.2^{0.75} = 43\times9.53 = \qty{410}{\celsius}$, saturated solidus $1331.1
- 410 = \qty{921}{\celsius}$. (b) The melt dissolves more water at higher pressure,
and the larger lowering cancels the \qty{118}{\celsius} rise of the dry solidus.
(c) The solidus must fall by $1331.1 - 1100 = \qty{231}{\celsius}$: $X =
(231/43)^{1/0.75} = 5.37^{1.333} = \qty{9.4}{\percent}$ of water, less than the
\qty{20.2}{\percent} at saturation, so the wedge can melt.
\end{solution}

\begin{solution}{pb:b2:mantle-melting:1}
\textbf{1.} $1350 + 273 = \qty{1623}{K}$.
\textbf{2.} $\num{4.0e-5}\times1623/(3300\times1000) = \qty{1.97e-8}{K/Pa} =
\qty{19.7}{K/GPa}$, rounded to \qty{20}{K/GPa}.
\textbf{3.} $\num{e9}/(3300\times9.81) = \qty{30.9}{km}$; $20/30.9 =
\qty{0.65}{K/km}$.
\textbf{4.} $1350 + 20P = 1085.7 + 132.9P - 5.1P^2$, that is $5.1P^2 - 112.9P +
264.3 = 0$: $P_0 = (112.9 - \sqrt{12746.4 - 5391.7})/10.2 = (112.9 - 85.76)/10.2 =
\qty{2.66}{GPa}$.
\textbf{5.} $2.66\times30.9 = \qty{82}{km}$; $T = 1350 + 20\times2.66 =
\qty{1403}{\celsius}$.
\textbf{6.} $1085.7 + 132.9\times2.66 - 5.1\times2.66^2 = 1085.7 + 353.5 - 36.1 =
\qty{1403}{\celsius}$.
\textbf{7.} $T_{\mathrm s} = 1085.7 + 199.4 - 11.5 = \qty{1273.6}{\celsius}$;
$T_{\mathrm l} = 1475.0 + 120.0 - 7.2 = \qty{1587.8}{\celsius}$; difference
\qty{314.2}{K}.
\textbf{8.} $1273.6 + 0.10\times314.2 = \qty{1305.0}{\celsius}$, \qty{1578}{K}.
\textbf{9.} $T_{\mathrm s}' = 132.9 - 15.3 = \qty{117.6}{K/GPa}$, $T_{\mathrm l}' =
80.0 - 9.6 = \qty{70.4}{K/GPa}$: $(\partial T/\partial P)_F = 117.6 +
0.10\times(70.4 - 117.6) = \qty{112.9}{K/GPa}$; $(\partial T/\partial F)_P =
\qty{314.2}{K}$.
\textbf{10.} $c_p/T = 1000/1578 = 0.634$; numerator $0.634\times112.9 - 12.1 = 71.6
- 12.1 = 59.5$; denominator $300 + 0.634\times314.2 = 499.2$: $0.119$, or
\qty{11.9}{\percent} per gigapascal.
\textbf{11.} With $F = T'^{1.5}$ the temperature must rise steeply with $F$ to make
the first per cent of melt, so $(\partial T/\partial F)_P$ is large near the solidus
and melting starts slowly. $F_{\max} = 0.09\times2.66 = 0.24$, \qty{24}{\percent}.
\textbf{12.} $0.13/0.50 = 0.26$: clinopyroxene runs out at \qty{26}{\percent} of
melting. At \qty{24}{\percent} the residue beneath the axis still holds a little:
a \omterm{def:b2:mantle-melting:peridotites}{lherzolite} poor in clinopyroxene, on the edge of \omterm{def:b2:mantle-melting:peridotites}{harzburgite}.
\textbf{13.} $h_c = 0.24\times82/2 = \qty{9.8}{km}$.
\textbf{14.} $F_{\max}/2 = \qty{12}{\percent}$.
\textbf{15.} \qty{9.8}{km} is \qty{38}{\percent} more than \qty{7.1}{km}, above the
range of normal crust.
\textbf{16.} Constant productivity gives the deepest mantle too much melt; the
fitted model makes the first melt slowly, $F(z)$ stays lower near the onset, and
$\int F\,\mathrm dz$ is smaller.
\textbf{17.} $F_{\max} = \Pi P_0$ and $H = 30.9P_0$ (km): $h_c = \Pi P_0\times30.9P_0/2
= kP_0^2$ with $k = 0.09\times30.9/2 = \qty{1.39}{km/GPa^2}$.
\textbf{18.} $P_0 = \sqrt{7.1/1.39} = \sqrt{5.11} = \qty{2.26}{GPa}$.
\textbf{19.} $2.26\times30.9 = \qty{70}{km}$.
\textbf{20.} $T_{\mathrm s}(2.26) = 1085.7 + 300.4 - 26.0 = \qty{1360.0}{\celsius}$,
and $T_p = 1360.0 - 20\times2.26 = 1360.0 - 45.2 = \qty{1315}{\celsius}$.
\textbf{21.} For \qty{6.3}{km}: $P_0 = \sqrt{4.53} = \qty{2.13}{GPa}$, $T_{\mathrm s}
= \qty{1345.6}{\celsius}$, $T_p = \qty{1303}{\celsius}$. For \qty{7.9}{km}: $P_0 =
\sqrt{5.68} = \qty{2.38}{GPa}$, $T_{\mathrm s} = \qty{1373.1}{\celsius}$, $T_p =
\qty{1326}{\celsius}$. The scatter of the crust allows $1315 \pm
\qty{12}{\celsius}$.
\textbf{22.} $P_0 = \sqrt{20/1.39} = \sqrt{14.39} = \qty{3.79}{GPa}$; $T_{\mathrm s} =
1085.7 + 503.7 - 73.3 = \qty{1516.1}{\celsius}$; $T_p = 1516.1 - 75.8 =
\qty{1440}{\celsius}$, an excess of $1440 - 1315 = \qty{125}{\celsius}$.
\textbf{23.} The two use different evidence. The lavas record the hottest mantle,
at the centre of the plume; the crust averages the melts of all the mantle that
passes through the melting region, much of it cooler. The simple model also keeps
the productivity constant, although clinopyroxene runs out at \qty{26}{\percent}
while $\Pi P_0 = 0.09\times3.79 = 0.34$ here: real hot mantle melts less, and needs
a higher temperature to make the same crust.
\textbf{24.} $T_{\mathrm s}'(2.26) = 132.9 - 23.1 = \qty{109.8}{K/GPa}$; $\Delta P_0 =
10/(109.8 - 20) = \qty{0.111}{GPa}$; since $h_c \propto P_0^2$, $\Delta h_c/h_c =
2\times0.111/2.26 = \qty{9.8}{\percent}$, about a tenth per \qty{10}{\celsius}.
\textbf{25.} At a very slow ridge the mantle rises so slowly that it loses heat by
conduction towards the cold surface; it leaves the adiabat, its melting stops
deeper, and the top of the triangle, where $F$ is largest, is cut off: the crust is
thinner.
\textbf{26.} The crust of normal ridges gives a \omterm{def:b2:mantle-melting:tp}{mantle potential temperature} of
\textbf{\qty[text-series-to-math=true]{1315}{\celsius}}, within about
\qty{12}{\celsius}: \qty{35}{\celsius} below the \qty{1350}{\celsius} of mineral
physics, which is uncertain by about as much, and well inside the range 1280 to
\qty{1400}{\celsius} from lavas. Three independent thermometers agree on a mantle
near \qty{1300}{\celsius} to \qty{1350}{\celsius}.
\end{solution}
